\documentclass[border=1pt]{standalone}

\usepackage[T1]{fontenc}
\usepackage[utf8]{inputenc}
\usepackage{lmodern}
\usepackage{microtype}
\usepackage{graphicx}
\usepackage{booktabs}
\usepackage{amsmath,amssymb,amsfonts}
\usepackage{bm}
\usepackage{siunitx}



\usepackage{tikz}
\usepackage{pgfplots}
\usepackage{pgfplotstable}
\usepackage{caption}
\pgfplotsset{compat=newest}




\newcommand{\papertitle}{Spectroscopic Constraints on Kelvin-Like Internal Modes in Effective Filament Models of Atomic Bound States}
\newcommand{\papershorttitle}{Kelvin-Like Internal Modes in Atomic Bound States}
\newcommand{\paperdoi}{10.5281/zenodo.18018084}
\newcommand{\paperorcid}{0009-0006-1686-3961}
\newcommand{\paperemail}{info@omariskandarani.com}
\newcommand{\papercity}{Groningen}
\newcommand{\papercountry}{The Netherlands}
\newcommand{\paperkeywords}{Kelvin modes, effective filament models, atomic spectroscopy, excitation gap, low-energy reduction}

\raggedbottom
\\pgfplotsset{compat=1.18}


\providecommand{\papertitle}{Spectroscopic Constraints on Kelvin-Like Internal Modes in Effective Filament Models of Atomic Bound States}
\providecommand{\papershorttitle}{Kelvin-Like Internal Modes in Atomic Bound States}
\providecommand{\paperdoi}{10.5281/zenodo.18018084}
\providecommand{\paperorcid}{0009-0006-1686-3961}
\providecommand{\paperemail}{info@omariskandarani.com}
\providecommand{\papercity}{Groningen}
\providecommand{\papercountry}{The Netherlands}
\providecommand{\paperkeywords}{Kelvin modes, effective filament models, atomic spectroscopy, excitation gap, low-energy reduction}
\providecommand{\epsilonK}{\varepsilon_K}

\usetikzlibrary{
    spath3,
    intersections,
    arrows,
    arrows.meta,
    knots,
    calc,
    hobby,
    decorations.pathreplacing,
    shapes.geometric,
}

% ---------------- Global styles (safe for 'knots') ----------------
\tikzset{
    knot diagram/every strand/.append style={
        line cap=round,
        line join=round,
        ultra thick,
        black
    },
    every knot/.style={line cap=round,line join=round,very thick},
    strand/.style={line cap=round,line join=round,line width=3pt,draw=black},
    over/.style={preaction={draw=white,line width=6.5pt}},
% DO NOT set a global "every path" here; it breaks internal clip paths.
}
\providecommand{\swirlarrow}{\rightsquigarrow}
\begin{document}
\begin{tikzpicture}
[
                width=12cm,
                height=8cm,
                xlabel={$T$ [K]},
                ylabel={$U(T)$ [arb. units]},
                legend style={at={(0.05,0.95)}, anchor=north west},
                domain=0:1.5,
                samples=100,
                axis lines=left,
                enlargelimits=true,
                grid=both,
                xtick={0.5,1.0,1.5},
                ytick={0,1,2,3,4},
                xlabel style={font=\small},
                ylabel style={font=\small},
                title={Kelvin-Mode Thermal Energy vs Temperature},
                title style={font=\bfseries\small},
                thick
                ]

                \addplot[blue,thick] {4*x^2};
                \addlegendentry{Ungapped case: illustrative algebraic example}

                \addplot[red,thick,dashed] {5*exp(-1/x)};
                \addlegendentry{Gapped case: exponential suppression}
                


                [shift={(7.8,4.3)}]
                    [
                    width=4.2cm,
                    height=3.5cm,
                    title={Kelvin Spectrum},
                    title style={font=\bfseries\tiny},
                    xlabel={$k$},
                    ylabel={$E(k)$},
                    axis lines=left,
                    ticks=none,
                    xmin=0, xmax=3.5,
                    ymin=0, ymax=4.5,
                    xtick=\empty,
                    ytick=\empty,
                    clip=false
                    ]
                    \fill[gray!20] (axis cs:0,0) rectangle (axis cs:3.5,1.5);
                    \draw[dashed, gray] (axis cs:0,1.5) -- (axis cs:3.5,1.5);
                    \foreach \x in {0.5,1.0,1.5,2.0,2.5,3.0} {
                        \addplot[only marks, mark=*, red] coordinates {(\x, {1.5 + 0.5*\x})};
                    }
                    \node[align=center,font=\tiny] at (axis cs:1.75,0.75) {No modes\\[-0.3em]($E<\Delta_K$)};
\end{tikzpicture}
\end{document}